\chapter*{Résumé}
\phantomsection 
Ce projet présente la conception et le développement de \textbf{BViral Control Center}, une plateforme web intelligente dédiée à la gestion de contenu sur les réseaux sociaux et au traitement vidéo assisté par l’intelligence artificielle. Elle permet aux créateurs et administrateurs de gérer plusieurs comptes, planifier et publier des contenus, suivre les statistiques, et recevoir des alertes via un tableau de bord centralisé.

La plateforme intègre également un \emph{AI Video Studio} capable d’améliorer la qualité des vidéos, de générer automatiquement des sous‑titres et de prédire la performance d’un contenu avant publication. Un pipeline multimodal a été mis en place pour extraire des caractéristiques visuelles, audio et textuelles à partir de vidéos courtes. Ces caractéristiques alimentent des modèles d’apprentissage automatique destinés à estimer la portée et l’engagement futurs.

L’architecture du système a été formalisée à l’aide de diagrammes UML (cas d’usage, classes, séquences) afin de définir les aspects fonctionnels et structurels de l’application. La solution combine des technologies web modernes, des services cloud et des outils d’IA avancés pour offrir un environnement évolutif et performant dédié à l’optimisation du contenu numérique.

